\begin{ouverture}[espace]{16}{Espace, solides et repérage sur la sphère}{Partie 4 · Géométrie}
{Perspectives · Patrons · Sections · Volumes · Latitude et longitude}
\objectif{Interpréter plusieurs représentations d’un solide.}
\objectif{Calculer des volumes et des assemblages.}
\objectif{Se repérer dans un pavé et sur une sphère.}
\prerequis{Aires usuelles, Pythagore et unités de volume.}
\end{ouverture}
\pagestyle{otMain}\essentiel\fontsize{11}{15}\selectfont
\essSection{1}{Représenter et mesurer les solides}
\begin{savoirfaire}[Lire une représentation]Une perspective cavalière représente les arêtes cachées en pointillés. Les longueurs dessinées ne sont pas toutes les longueurs réelles. Une vue de face ou de dessus projette les formes ; une section est l’intersection avec un plan. Un patron permet de replier les faces sans chevauchement.\end{savoirfaire}
\begin{propriete}Pour un pavé : $V=L\ell h$. Pour un prisme droit ou un cylindre : $V=\mathcal A_{\rm base}h$. Pour une pyramide ou un cône : $V=\mathcal A_{\rm base}h/3$. La hauteur est perpendiculaire à la base.\end{propriete}
\begin{propriete}Pour une boule de rayon r : $V=4\pi r^3/3$. Pour sa surface, la sphère : $\mathcal A=4\pi r^2$. Un diamètre vaut deux rayons.\end{propriete}
\begin{savoirfaire}[Assembler]Décomposer en solides usuels, additionner des volumes sans recouvrement ou soustraire un creux. Conserver $\pi$ jusqu’au résultat final. Pour un réservoir, convertir ensuite en litres.\end{savoirfaire}
\clearpage
\essSection{2}{Sections et repérages}
\begin{propriete}La section d’un pavé par un plan parallèle à une face est un rectangle de mêmes dimensions que cette face. La section d’une boule par un plan est un disque ; son contour sur la sphère est un cercle. Le cercle est maximal quand le plan passe par le centre.\end{propriete}
\begin{savoirfaire}[Un repère dans un pavé]Fixer une origine et trois directions d’arêtes perpendiculaires. Les coordonnées $(x;y;z)$ donnent trois déplacements. Préciser l’unité et les limites de chaque coordonnée.\end{savoirfaire}
\begin{definition}Sur une sphère terrestre, la latitude indique un angle au nord ou au sud de l’équateur, de 0 à 90 degrés. La longitude indique un angle à l’est ou à l’ouest du méridien d’origine, de 0 à 180 degrés. Les pôles ont une latitude de 90 degrés ; leur longitude n’est pas unique.\end{definition}
\begin{attention}Deux points de même latitude appartiennent au même parallèle ; de même longitude, au même méridien. Une carte plane peut déformer distances et aires. Les exercices géographiques donnent des positions fictives ou un modèle explicite.\end{attention}
\clearpage\rubrique{Découverte}
\begin{blocLivre}{ex-16-1}\begin{exercice}[1]{Un pavé}
Un pavé a pour dimensions 8 cm, 5 cm et 3 cm. Calculer son volume.
\label{fig-ex-16-1}\figureGeo{espace-01}
\end{exercice}\end{blocLivre}
\begin{blocLivre}{ex-16-2}\begin{exercice}[1]{Un cylindre}
Un cylindre a un rayon de 3 cm et une hauteur de 10 cm. Donner son volume exact et au cm$^3$.
\label{fig-ex-16-2}\figureGeo{espace-02}
\end{exercice}\end{blocLivre}
\begin{blocLivre}{ex-16-3}\begin{exercice}[1]{Une pyramide}
Une pyramide a une base carrée de côté 6 cm et une hauteur de 9 cm. Calculer son volume.
\label{fig-ex-16-3}\figureGeo{espace-03}
\end{exercice}\end{blocLivre}
\begin{blocLivre}{ex-16-4}\begin{exercice}[1]{Un cône}
Un cône a pour diamètre 8 cm et hauteur 9 cm. Calculer son volume exact.
\label{fig-ex-16-4}\figureGeo{espace-04}
\end{exercice}\end{blocLivre}
\begin{blocLivre}{ex-16-5}\begin{exercice}[1]{Une boule}
Une boule a un rayon de 3 cm. Donner son volume et l’aire de sa sphère, sous forme exacte.
\label{fig-ex-16-5}\figureGeo{espace-05}
\end{exercice}\end{blocLivre}

\clearpage\rubrique{Entraînement}
\begin{blocLivre}{ex-16-6}\begin{exercice}[2]{Lire un patron de cube}
Le patron comporte six carrés de côté 2 cm. Nommer les faces opposées à A, B et C après pliage ; calculer l’aire totale du cube.
\label{fig-ex-16-6}\figureGeo{espace-patron}
\end{exercice}\end{blocLivre}
\begin{blocLivre}{ex-16-7}\begin{exercice}[2]{Trois vues}
Un pavé a une longueur de 8 cm, une profondeur de 5 cm et une hauteur de 3 cm. Donner les dimensions des vues de face, de dessus et du côté droit.
\label{fig-ex-16-7}\figureGeo{espace-vues}
\end{exercice}\end{blocLivre}
\begin{blocLivre}{ex-16-8}\begin{exercice}[2]{Une section de pavé}
Un pavé de base 8 cm sur 5 cm et de hauteur 6 cm est coupé par un plan parallèle à la base à 2 cm de celle-ci. Quelle forme et quelle aire a la section ?
\label{fig-ex-16-8}\figureGeo{espace-section-pave}
\end{exercice}\end{blocLivre}
\begin{blocLivre}{ex-16-9}\begin{exercice}[2]{Des coordonnées dans un pavé}
Dans le pavé de dimensions 6, 4, 3 unités, O est l’origine ; les axes suivent longueur, profondeur, hauteur. Donner le sommet opposé à O et la projection de P(2;3;1) sur le plan de base.
\label{fig-ex-16-9}\figureGeo{espace-repere}
\end{exercice}\end{blocLivre}
\begin{blocLivre}{ex-16-10}\begin{exercice}[2]{Latitude et longitude}
Les points fictifs A et B ont pour positions $(35^\circ\text{ N};20^\circ\text{ E})$ et $(35^\circ\text{ N};50^\circ\text{ O})$. Sont-ils sur le même parallèle ? Le même méridien ? Quelle est la différence de longitude ?
\label{fig-ex-16-10}\figureGeo{espace-sphere}
\end{exercice}\end{blocLivre}

\clearpage\rubrique{Problèmes}
\begin{blocLivre}{ex-16-11}\begin{exercice}[3]{Un réservoir assemblé}
Un réservoir fermé est formé d’un cylindre de rayon 1 m et de hauteur 2 m, surmonté d’une demi-boule de même rayon. Calculer sa capacité au litre.
\label{fig-ex-16-11}\figureGeo{espace-reservoir}
\end{exercice}\end{blocLivre}
\begin{blocLivre}{ex-16-12}\begin{exercice}[3]{Une boule creusée}
Une boule pleine de rayon 5 cm contient une cavité sphérique de rayon 2 cm, entièrement intérieure. Quel volume de matière reste-t-il ?
\label{fig-ex-16-12}\figureGeo{espace-cavite}
\end{exercice}\end{blocLivre}
\begin{blocLivre}{ex-16-13}\begin{exercice}[3]{Une section de sphère}
Une sphère de centre O a un rayon de 5 cm. Un plan la coupe à une distance OH = 3 cm du centre ; H est le centre du cercle de section. Calculer le rayon de ce cercle et l’aire du disque.
\label{fig-ex-16-13}\figureGeo{espace-section-sphere}
\end{exercice}\end{blocLivre}
\begin{blocLivre}{ex-16-14}\begin{exercice}[3]{Un cône réduit}
Un cône de rayon 6 cm et hauteur 12 cm est coupé parallèlement à sa base à mi-hauteur. Le petit cône de même sommet a un rapport de réduction $1/2$. Calculer son volume et celui du tronc restant.
\label{fig-ex-16-14}\figureGeo{espace-cone-coupe}
\end{exercice}\end{blocLivre}
\begin{blocLivre}{ex-16-15}\begin{exercice}[3]{La diagonale de l’espace}
Un pavé mesure 3 cm, 4 cm et 12 cm. Calculer une diagonale de la base 3 par 4 puis une diagonale intérieure joignant deux sommets opposés.
\label{fig-ex-16-15}\figureGeo{espace-diagonale}
\end{exercice}\end{blocLivre}

\clearpage
